% ============================================================================
% OPERATIONS RUNBOOK
% ============================================================================

\section{Operations Runbook}
\label{section:operations-runbook}

\subsection{Scope}

This section is the in-the-field guide. It assumes a Pi 5 has already
been provisioned per \texttt{docs/PI\_DEPLOYMENT.md} (one-time
hardware-and-OS setup) and that the topology is the one drawn in
\texttt{docs/ARCHITECTURE.md} (cockpit, Caddy, supervisor mux, k3s).
\textbf{This file answers only one question}: what does the operator
do during a session, when something is on fire, or when the routine
"power-on, drive, power-off" sequence drifts. Setup, build, and
topology are explicitly out of scope.

Throughout the runbook, \texttt{\$\{STAR\_ROOT\}} is whatever
absolute path the repo lives at on the Pi 5
(\texttt{/workspaces/STAR} on the development Pi; the value is
interchangeable). All paths in this section use that placeholder.

\subsection{Daily start sequence}

\textbf{Power up.} Plug the chassis battery, give it ~10 seconds for
the Pi 5 to come out of POST and finish boot. You should see the Pi
5's PWR LED solid red and the activity LED blinking; the RX72N PCB's
3.3 V LED solid green; the RPLiDAR LED solid green (the rotor will
spin up only when SLAM activates it).

\textbf{Wait for Pi services.} The Pi runs three units that come up
on boot, before any operator action:

\begin{itemize}
  \item \texttt{star-cpu-governor.service} (one-shot, pins governor
        to \texttt{schedutil}).
  \item \texttt{caddy.service} (TLS termination on
        \texttt{https://star.local}, binds to the tailscale interface
        \texttt{100.64.0.6}).
  \item \texttt{star-cockpit-api.service} (HTTP API at
        \texttt{127.0.0.1:9102} that the Grafana cockpit toggles
        autonomy / e-stop / speed against).
\end{itemize}

These are \emph{infrastructure} -- they always run. They do not
launch ROS2. To confirm the Pi is ready before you launch the stack:

\begin{quote}\small
\texttt{systemctl is-active caddy star-cockpit-api star-cpu-governor}
\end{quote}

All three should print \texttt{active}.

\textbf{Run \texttt{./start.sh}.} From \texttt{\$\{STAR\_ROOT\}} on
the Pi (over SSH or local console):

\begin{quote}\small
\texttt{./start.sh}
\end{quote}

\texttt{start.sh} is the canonical robot launcher. It:

\begin{enumerate}
  \item Calls \texttt{require\_host\_setup} -- warns (but does not
        block) if the Pi 5 host setup from
        \texttt{scripts/setup-pi5-host.sh} is incomplete.
  \item Brings up the WiFi AP \texttt{STAR-Robot} on
        \texttt{wlan0} unless \texttt{--no-ap} is passed. The Pi
        becomes \texttt{192.168.50.1}.
  \item Stops anything already running (\texttt{./stop.sh}), flushes
        the Cyclone DDS daemon's stale participant cache.
  \item Probes for an RX72N over SPI for ~4 seconds. Falls back to
        \texttt{virtual\_rx72n} (DEV mode) if the probe fails.
  \item Detects RPLiDAR on \texttt{/dev/rplidar} or
        \texttt{/dev/ttyUSB0}; stops \texttt{ModemManager} so it does
        not grab the port.
  \item Launches in order: \texttt{virtual\_rx72n} (DEV only),
        \texttt{star-gateway} (gRPC :50051, WS :8080),
        \texttt{star\_spi\_bridge}, fake odom TF (DEV only), SLAM
        stack (\texttt{slam.launch.py}), compliance engine,
        \texttt{star\_gateway\_bridge\_main}, UI (Vite, :5173), RViz
        (optional), Lichtblick web (:8081).
\end{enumerate}

The operator-facing flags are summarized below.

\begin{center}
\begin{tabular}{|p{3.3cm}|p{11.5cm}|}
\hline
\textbf{Flag} & \textbf{Meaning} \\
\hline
(no flags) & Auto-detect everything; full stack. Default field choice. \\
\hline
\texttt{--no-lidar} & Skip LiDAR / SLAM / EKF even when \texttt{/dev/rplidar} is present. Use for indoor bench tests where the rotor is in the way. \\
\hline
\texttt{--no-ui} & Skip the Vite dev server. Use when no operator laptop is connected; saves ~150 MB RAM. \\
\hline
\texttt{--no-compliance} & Skip the seven compliance nodes (ramp slope, doorway, protrusions, etc.). Use during pure-mobility testing. \\
\hline
\texttt{--no-ap} & Do not broadcast \texttt{STAR-Robot} AP. Use when the Pi is connected via Ethernet and you want to keep \texttt{wlan0} on a station network. \\
\hline
\texttt{--rviz} & Launch \texttt{rviz2} with \texttt{slam\_lidar.rviz} after startup. \\
\hline
\texttt{--help} & Print usage. \\
\hline
\end{tabular}
\end{center}

Environment overrides:
\texttt{STAR\_SIMULATION\_MODE=true|false} forces DEV / HW mode and
skips the SPI probe; \texttt{STAR\_WIFI\_SSID},
\texttt{STAR\_WIFI\_PASS}, and \texttt{STAR\_WIFI\_IFACE} override the
AP defaults.

\textbf{Verify health.} After \texttt{start.sh} prints its summary
banner, sanity-check the stack:

\begin{center}
\begin{tabular}{|p{4.4cm}|p{10.4cm}|}
\hline
\textbf{Check} & \textbf{Expected} \\
\hline
\texttt{nc -z localhost 50051} & gateway gRPC reachable. \\
\hline
\texttt{nc -z localhost 8080}  & gateway WS reachable. \\
\hline
\texttt{nc -z localhost 8765}  & \texttt{foxglove\_bridge} listening. \\
\hline
\texttt{nc -z localhost 5173}  & UI dev server reachable (HW WiFi mode prints \texttt{http://192.168.50.1:5173}). \\
\hline
\texttt{ros2 topic hz /scan}   & ~10\,Hz when LiDAR is up. \\
\hline
\texttt{ros2 topic hz /odom/unfiltered} & ~50\,Hz from \texttt{simple\_bridge} or \texttt{star\_gateway\_bridge}. \\
\hline
\texttt{ros2 topic echo /map -n 1 --field info} & a non-empty origin / resolution after the robot has moved. \\
\hline
\texttt{curl -s https://star.local/api/state} & \texttt{\{"autonomy":..., "estop":..., "speed":...\}} (cockpit API live). \\
\hline
\end{tabular}
\end{center}

\subsection{Daily stop sequence}

From \texttt{\$\{STAR\_ROOT\}}:

\begin{quote}\small
\texttt{./stop.sh}\\
\texttt{./stop.sh --stop-ap}      \quad (also tear down the WiFi AP)
\end{quote}

\texttt{stop.sh} unwinds in the reverse of startup, with
process-group SIGTERM first, a 3--5 second grace, then SIGKILL on
survivors:

\begin{enumerate}
  \item ROS2 client first: \texttt{star\_gateway\_bridge\_main}.
  \item ROS2 launch groups (\texttt{slam.launch.py},
        \texttt{compliance.launch.py},
        \texttt{star\_spi\_bridge.launch.py},
        \texttt{stereo\_camera.launch.py}).
  \item Individual ROS2 nodes that may have outlived their launch
        parent (\texttt{sllidar}, \texttt{slam\_toolbox},
        \texttt{ekf\_node}, \texttt{robot\_state\_publisher},
        \texttt{static\_transform\_publisher}, \texttt{gscam},
        \texttt{foxglove\_bridge}).
  \item Go binaries: \texttt{star-gateway}, \texttt{virtual\_rx72n}
        (a 2-second sleep here lets the serial port release before
        the next \texttt{start.sh}).
  \item Visualization: Vite UI, RViz, Lichtblick.
  \item A final pgrep sweep for orphans.
  \item By default, leaves the WiFi AP up so SSH from the operator
        laptop survives. \texttt{--stop-ap} tears it down.
\end{enumerate}

\textbf{When to use Ctrl-C vs \texttt{stop.sh} vs power-cycle.}
Ctrl-C in \texttt{start.sh}'s shell triggers the script's INT trap,
which calls \texttt{./stop.sh} -- equivalent. Use Ctrl-C if you ran
\texttt{start.sh} in the foreground and want a clean unwind. Use
\texttt{./stop.sh} from any other shell. Use a power-cycle only when
\texttt{stop.sh} fails to release \texttt{/dev/ttyACM*} or
\texttt{/dev/ttyUSB0} (rare; usually means a node is stuck in an
uninterruptible kernel I/O wait, which only a hardware reset clears).

\subsection{Dev mode}

\texttt{./dev.sh} is a one-liner wrapper:

\begin{quote}\small
\texttt{exec env STAR\_SIMULATION\_MODE=true \$\{STAR\_ROOT\}/start.sh --rviz "\$@"}
\end{quote}

It forces the SPI probe to skip, brings up \texttt{virtual\_rx72n}
(the Go simulator that emits synthesized telemetry), launches RViz,
and otherwise runs the same launch order as \texttt{start.sh}. Use it
when:

\begin{itemize}
  \item The robot hardware is not connected (desk-side dev).
  \item You want to test a new SLAM tuning on a recorded bag without
        the real motors.
  \item You want to debug the gateway or UI without an RX72N
        responding.
\end{itemize}

DEV mode publishes a static \texttt{odom -> base\_link} TF
(\texttt{robot stationary at origin}) so the SLAM TF chain stays
intact. Never \texttt{./dev.sh} on a robot whose wheels are on the
ground.

\subsection{Field deployment checklist}

\begin{center}
\begin{tabular}{|p{2.6cm}|p{4.0cm}|p{8.2cm}|}
\hline
\textbf{Phase} & \textbf{Item} & \textbf{Pass criterion} \\
\hline
Pre-flight & Power cable & Battery shows $>$11.5 V at the chassis
distribution block. \\
\hline
Pre-flight & USB CDC enumerated &
\texttt{ls /dev/star-rx72n-proto /dev/ttyACM*} resolves; if
missing, see the E2 Lite recovery path inside
\texttt{scripts/slam-mvp.sh}'s \texttt{ensure\_mcu()}. \\
\hline
Pre-flight & RPLiDAR powered &
\texttt{ls /dev/rplidar /dev/ttyUSB0}; LED on the lidar should be
solid green. \\
\hline
Pre-flight & IMU calibration & BNO055 reaches sys / gyro / accel
calib level 3 -- check via the cockpit IMU panel or
\texttt{ros2 topic echo /star/imu/calib\_status -n 1}. \\
\hline
Pre-flight & Host setup &
\texttt{./start.sh} prints no \texttt{require\_host\_setup}
warnings. If it does, run
\texttt{sudo bash \$\{STAR\_ROOT\}/scripts/setup-pi5-host.sh}. \\
\hline
Operational & E-stop test &
\texttt{ros2 topic pub --once /star/estop std\_msgs/Bool '\{data: true\}'}
-- wheels must brake within $\le$50\,ms. Confirm with cockpit /
Grafana. \\
\hline
Operational & Autonomy toggle &
\texttt{POST https://star.local/api/autonomy/toggle} -- the cockpit
or \texttt{ros2 topic echo /star/state/autonomy\_enable} should
flip. \\
\hline
Operational & Supervisor mux &
With autonomy off, \texttt{/cmd\_vel} drives the robot; with
autonomy on, only \texttt{/nav2/cmd\_vel} drives. Use
\texttt{ros2 topic echo /cmd\_vel\_out}. \\
\hline
Post-flight & Log archive &
\texttt{tar czf run-\$(date +\%Y\%m\%d-\%H\%M).tgz /tmp/star-logs/}
-- copy off the Pi before shutdown. \\
\hline
Post-flight & Mission report &
Compliance engine writes a PDF under
\texttt{\$\{STAR\_ROOT\}/star-compliance/reports/} -- archive
alongside the logs. \\
\hline
\end{tabular}
\end{center}

\subsection{Cockpit URLs and ports}

These are the ports operators interact with. All bind to
\texttt{0.0.0.0} on the Pi 5 unless noted otherwise (so they are
reachable from a laptop joined to \texttt{STAR-Robot} or over the
tailnet).

\begin{center}
\begin{tabular}{|p{1.5cm}|p{4.4cm}|p{8.9cm}|}
\hline
\textbf{Port} & \textbf{Service} & \textbf{What it serves; credential strategy} \\
\hline
443  & Caddy on Pi 5 & TLS terminator. Bound to
\texttt{100.64.0.6} (tailscale interface) only.
Routes \texttt{/api/*} to the cockpit API at
\texttt{100.64.0.6:9102}; everything else proxies to
\texttt{star\_simple\_bridge} on \texttt{localhost:9101}. Internal
CA -- one-time install on operator laptop. \\
\hline
50051 & gateway gRPC & Used by the ROS2
\texttt{star\_gateway\_bridge} only. Localhost-only by convention;
no auth. \\
\hline
8080  & gateway WS & UI -> gateway WebSocket.
\texttt{WS\_STRICT\_ORIGIN=false} during dev. No auth (lab-only). \\
\hline
8765  & \texttt{foxglove\_bridge} & ROS topic visualizer used by
both the bundled Lichtblick at \texttt{:8081} and any external
Foxglove Studio install (\texttt{ws://<pi5>:8765}). Plaintext;
relies on tailscale or the AP for confidentiality. \\
\hline
5173  & UI Vite dev server & The TypeScript operator UI. No auth,
plaintext. AP-network only by default. \\
\hline
8081  & Lichtblick web (\texttt{npx serve}) & Self-hosted Foxglove
Studio fork. Requires one-time install via
\texttt{scripts/setup-lichtblick-web.sh}. \\
\hline
9100  & node\_exporter & Prometheus host metrics scraped by Alloy
on the Pi and pushed to k3s. No auth. \\
\hline
9101  & \texttt{star\_simple\_bridge} & Serves
\texttt{/map.\{png,pgm,yaml\}}, accepts
\texttt{POST /map/reset?confirm=yes}, exposes its own custom
metrics. Fronted by Caddy on \texttt{star.local}. \\
\hline
9102  & \texttt{star\_cockpit\_api} & The Python HTTP backend for
the cockpit toggles. Binds to \texttt{100.64.0.6:9102}; reachable
only via Caddy at \texttt{https://star.local/api/*}. Endpoints:
\texttt{GET /api/state}, \texttt{POST /api/autonomy/toggle},
\texttt{POST /api/estop/toggle},
\texttt{POST /api/slam/\{start,stop\}},
\texttt{POST /api/cmd\_vel}, \texttt{POST /api/speed}. \\
\hline
\end{tabular}
\end{center}

The Grafana cockpit dashboard (panels: SLAM map, Motors / FL.. BR
velocity / duty / faults, IMU roll / pitch / heading / gyro / accel,
Odometry linear / angular / pose, LiDAR closest obstacle / bearing /
range stats, Bridge health line-rate / parse errors / serial reopens)
runs in the homelab k3s cluster, not on the Pi. It scrapes via Alloy
and embeds \texttt{https://star.local/map.png} directly from the
operator's browser.

\subsection{Systemd unit reference}

\begin{center}
\begin{tabular}{|p{4.4cm}|p{4.0cm}|p{6.4cm}|}
\hline
\textbf{Unit} & \textbf{What it runs} & \textbf{Status / log access} \\
\hline
\texttt{caddy.service} & \texttt{caddy run --config /etc/caddy/Caddyfile} as user \texttt{caddy}. Provides TLS on \texttt{star.local}, binds to \texttt{100.64.0.6}. \texttt{Restart=always}. & \texttt{systemctl status caddy};
\texttt{journalctl -u caddy -f}. \\
\hline
\texttt{star-cockpit-api.service} & Sources ROS2, exports
\texttt{rmw\_cyclonedds\_cpp} + \texttt{CYCLONEDDS\_URI}, runs
\texttt{python3 /opt/star\_cockpit\_api/server.py}. Subscribes
to latched supervisor state topics, publishes to
\texttt{/star/autonomy\_enable}, \texttt{/star/estop},
\texttt{/star/speed\_scale}. \texttt{Restart=always},
\texttt{RestartSec=3}. & \texttt{systemctl status star-cockpit-api};
\texttt{journalctl -u star-cockpit-api -f}. \\
\hline
\texttt{star-slam-mvp.service} & \emph{Optional, opt-in.} Runs
\texttt{scripts/slam-mvp.sh start} as user \texttt{star} after
\texttt{network-online.target} and a 8-second sleep so USB
devices have enumerated. \texttt{Type=forking,
RemainAfterExit=yes, Restart=no}. \textbf{Not enabled by
default.} Enable with \texttt{sudo systemctl enable --now
star-slam-mvp.service}. & \texttt{systemctl status star-slam-mvp};
log files live under \texttt{/tmp/slam-mvp-logs/<node>.log}. \\
\hline
\texttt{star-cpu-governor.service} & One-shot at very early boot
(\texttt{DefaultDependencies=no, Before=base.target}); writes
\texttt{schedutil} into every CPU's \texttt{scaling\_governor}.
\texttt{RemainAfterExit=yes} so \texttt{systemctl status} reflects
"applied". & \texttt{systemctl status star-cpu-governor}; not in
journal beyond the initial run. \\
\hline
\texttt{star-robot.service} & Optional auto-boot wrapper around
\texttt{./start.sh --no-ui} / \texttt{./stop.sh}. Installed by
\texttt{setup-pi5-host.sh} but \textbf{disabled} by default --
auto-starting on a robot the operator did not intend to drive is
a safety hazard. & Enabled only for field deployments;
\texttt{journalctl -u star-robot} captures
\texttt{/tmp/star-logs/systemd.log}. \\
\hline
\end{tabular}
\end{center}

The full safety model behind autonomy / e-stop arbitration is in
\S\ref{section:safety-architecture}; the supervisor topic mux table
in \texttt{docs/ARCHITECTURE.md} is the authoritative version.

\subsection{Wi-Fi and Ethernet helpers}

Two helper scripts live on the Pi 5 \texttt{PATH} (installed by
\texttt{setup-pi5-host.sh}, sourced from
\texttt{\$\{STAR\_ROOT\}/scripts/}):

\begin{center}
\begin{tabular}{|p{3.0cm}|p{11.7cm}|}
\hline
\textbf{Helper} & \textbf{Behaviour} \\
\hline
\texttt{star-ap up} & Brings up the \texttt{STAR-Hotspot}
NetworkManager profile on \texttt{wlan0}. Pi becomes
\texttt{192.168.50.1} broadcasting SSID \texttt{STAR-Robot}.
\textbf{Drops any existing wlan0 SSH session.} \\
\hline
\texttt{star-ap down} & Tears down the AP and lets
NetworkManager rejoin the previous station network. \\
\hline
\texttt{star-ap status} & Reports current Wi-Fi mode (AP / station / down). \\
\hline
\texttt{star-eth up} & Force-activates the
\texttt{netplan-eth0} shared profile. Pi becomes
\texttt{10.42.0.1}, runs dnsmasq for the laptop on the other end.
The autoconnect default is on, so this is rarely needed -- plug in
the cable and \texttt{star-eth status} should already report
\texttt{UP}. \\
\hline
\texttt{star-eth down} & Drops the eth0 shared profile (cable can
stay plugged in). \\
\hline
\texttt{star-eth auto on|off} & Toggles
\texttt{connection.autoconnect} on the profile.
\texttt{auto off} means the cable does nothing until you call
\texttt{star-eth up}. \\
\hline
\texttt{star-eth status} & Shows carrier, IPv4, peer DHCP lease,
and the autoconnect setting. \\
\hline
\end{tabular}
\end{center}

The two are mutually exclusive in practice -- \texttt{star-ap up} owns
\texttt{wlan0}; \texttt{star-eth} owns \texttt{eth0}. Running both is
fine; the Pi just routes per-interface. SSH options are summarized
in the cheat sheet at \texttt{scripts/star-help.sh} (run
\texttt{star-help} on the Pi for a colorized rendering).

\subsection{Common runtime issues}

For the deep gotcha catalog -- USB enumeration, AD2 buffer caps, FPU
context save, MSTPCR module-stop -- see
\S\ref{section:troubleshooting}. The few entries below are
operations-specific and not duplicated there.

\paragraph{Robot moves but no telemetry on the cockpit / UI.}
\texttt{/cmd\_vel\_out} is being consumed (motors spin) but the UI
shows stale topics. Almost always the WebSocket origin or URL: the
gateway runs with \texttt{WS\_STRICT\_ORIGIN=false} during dev;
verify with \texttt{ss -tlnp | grep 8080} that the gateway is on
\texttt{:8080}, then check the UI is pointing at
\texttt{ws://192.168.50.1:8080/ws} (AP) or
\texttt{ws://10.42.0.1:8080/ws} (Eth). If the cockpit (Grafana
panels) has no telemetry but Lichtblick over
\texttt{ws://<pi>:8765} is fine, the bridge to k3s Prometheus is the
issue, not the robot.

\paragraph{Robot freezes mid-motion.}
The RX72N IWDT will reset the chip if \texttt{watchdog\_task} ever
stops kicking, which forces motor outputs to zero on power-up. See
\S\ref{section:safety-architecture} for the three layers
(supervisor / comm / firmware-watchdog) and which one likely caught
the fault. Recovery: power-cycle the chassis. The Pi 5 stays up; only
the RX72N reboots.

\paragraph{Compass / yaw is bad after start.}
BNO055 needs the calibration loop to converge. Wave the chassis in a
figure-eight for 15--30 seconds; sys / gyro / accel calib levels
should each climb to 3. Until they do, the EKF heading drifts and
SLAM matches poorly. The cockpit IMU panel exposes the calib state
directly; \texttt{ros2 topic echo /star/imu/calib\_status -n 1} is
the CLI equivalent.

\paragraph{LiDAR reports zero ranges across the whole scan.}
Power, not data. Either:

\begin{itemize}
  \item The RPLiDAR's 5\,V rail dropped (battery low; chassis
        distribution loose) -- LED on the lidar will be off or
        amber.
  \item \texttt{ModemManager} grabbed \texttt{/dev/ttyUSB0} during
        boot before \texttt{start.sh} could stop it. Run
        \texttt{sudo systemctl stop ModemManager} and
        \texttt{./start.sh} again.
  \item Permissions: confirm \texttt{ls -l /dev/ttyUSB0} shows
        group-writable; the udev rule from
        \texttt{setup-pi5-host.sh} should have set this.
\end{itemize}

\subsection{Log locations}

\begin{center}
\begin{tabular}{|p{4.6cm}|p{4.4cm}|p{5.7cm}|}
\hline
\textbf{Source} & \textbf{Path} & \textbf{Access pattern} \\
\hline
\texttt{start.sh} children &
\texttt{/tmp/star-logs/<node>.log} &
\texttt{tail -f /tmp/star-logs/gateway.log}, etc. Rotated daily by
\texttt{/etc/logrotate.d/star-robot}, 3 generations, 20\,MB cap. \\
\hline
\texttt{slam-mvp.sh} children &
\texttt{/tmp/slam-mvp-logs/<node>.log} &
\texttt{scripts/slam-mvp.sh logs <node>} or
\texttt{tail -f /tmp/slam-mvp-logs/<node>.log}. \\
\hline
\texttt{star-cockpit-api.service} & journald only &
\texttt{journalctl -u star-cockpit-api -f}; backs to
\texttt{/var/log/journal/}. \\
\hline
\texttt{caddy.service} & journald only &
\texttt{journalctl -u caddy -f}. \\
\hline
\texttt{star-slam-mvp.service} & journald + per-node files &
\texttt{journalctl -u star-slam-mvp} for the systemd wrapper;
per-node detail under \texttt{/tmp/slam-mvp-logs/}. \\
\hline
\texttt{star-robot.service} (if enabled) &
\texttt{/tmp/star-logs/systemd.log} +
\texttt{journalctl -u star-robot} & As specified in
\texttt{StandardOutput=append:} of the unit. \\
\hline
\texttt{rosout} log files & \texttt{\textasciitilde/.ros/log/<run-id>/} &
ROS2 daemon writes one directory per run; useful for crash
post-mortems. \\
\hline
\texttt{nginx-style HTTP access} (Caddy) & journald only &
\texttt{journalctl -u caddy | grep "GET /api"} for cockpit-API
requests. \\
\hline
\end{tabular}
\end{center}

To grep across all robot logs from one shell:

\begin{quote}\small
\texttt{grep -RIn ERROR /tmp/star-logs/ /tmp/slam-mvp-logs/}
\end{quote}

\subsection{Ops references}

The placeholder \texttt{\$\{STAR\_ROOT\}} is whatever absolute path
the repository lives at on the host -- \texttt{/workspaces/STAR} on
the development Pi 5, \texttt{\textasciitilde/STAR} on a fresh
checkout, etc. Substitute as appropriate; the systemd units in
\texttt{infra/pi5/etc/systemd/system/} hardcode
\texttt{/workspaces/STAR} and must be re-templated if the repo lives
elsewhere.

For a deeper map of every directory referenced in this section
(\texttt{infra/pi5/}, \texttt{scripts/}, \texttt{star-ros2/launch/},
the systemd units, \texttt{monitoring/}), see
\S\ref{section:repository-layout}. For setup steps that run once per
Pi (deploy keys, apt packages, host-level udev / governor), see
\texttt{docs/PI\_DEPLOYMENT.md}. For the supervisor mux topology and
TLS / hostname strategy, see \texttt{docs/ARCHITECTURE.md}.
